\section{How the 2026 figures are produced}\label{sec:upmodel}

The construction-era tables (\appref{app:construction}) were generated from a
backing spreadsheet. The 2026 figures come from a sizing model maintained as code in
\href{https://github.com/lsst/rubin-sizing-model/tree/\smSourceRef}{\texttt{lsst/rubin-sizing-model}}.
Every 2026 table and model figure was generated from version
\texttt{\smSourceRef} of that repository; none were typed in by hand.

One parameter file, \texttt{sizing\_params.yaml}, holds every input number,
label and note. A script, \texttt{generate\_model.py}, turns it into an Excel
workbook of live formulas,
\begin{center}
\small\href{https://github.com/lsst/rubin-sizing-model/blob/\smSourceRef/rubin_usdf_model_2026_condensed_with_pricing.xlsx}{\texttt{rubin\_usdf\_model\_2026\_condensed\_with\_pricing.xlsx}}
\end{center}
The workbook has two \emph{reference} tabs holding every raw input: sizing
inputs, and costs with the purchase history. Six \emph{calculated} tabs
derive everything else from them:
\begin{itemize}
\item all processing at the USDF;
\item the facility split;
\item Qserv;
\item a blank calculator for partner facilities;
\item two USDF cost forecasts, one with all processing at the USDF and one
      with DRP shared.
\end{itemize}
Yellow cells mark the assumptions that most move the answer.

\subsection{Regenerating the tables}\label{sec:upregen}

The 2026 tables, and the macros that supply the figures quoted in
the text, are written by the same script. To reproduce them from the version
cited here:
\begin{center}
\begin{minipage}{0.92\textwidth}
\small
\begin{alltt}
git clone --branch \smSourceRef{} https://github.com/lsst/rubin-sizing-model
cd rubin-sizing-model
uv run generate_model.py --emit-tex ../dmtn-135/update2026/tables
\end{alltt}
\end{minipage}
\end{center}
This needs \texttt{git} and \href{https://docs.astral.sh/uv/}{\texttt{uv}}.

The model versions cited are listed with their tags in
\texttt{sizing-models.toml}; \texttt{make sizing-tables} regenerates them, as
the CI build does. The construction-era tables still come from the Google
backing sheet (\texttt{make tables}).

\subsection{Sources}

\begin{itemize}
\item \href{https://rubinobs.atlassian.net/wiki/spaces/DM/pages/1062404112/DRP+Resource+Usage+Estimates}{\emph{DRP Resource Usage Estimates}}
      (Confluence, last modified 2026-03-23), the source for DRP compute and
      storage, using its most recent pipeline version:
      \begin{itemize}
      \item DR1 compute, \S1.2.4: runtimes from \vthirty{} (DM-53697) on
            Milano nodes, for the first year of \texttt{baseline\_v5.0.0}
            (\S1.2.5). This supersedes the \wfortyone{} estimate in \S1.2.3.
      \item DR1 storage by dataset type with compression, \S1.2.2, using
            dataset sizes from the same run.
      \item DP2 storage with compression, \S1.1.3, for the $\sim$30k-visit DP2
            visit list.
      \end{itemize}
      \href{https://rtn-109.lsst.io/}{RTN-109} describes the method; its
      figures predate compression and v30.
\item \emph{Rubin Key Numbers} (Confluence): observing nights, visits per
      night, image size and raw compression.
\item \emph{DP2 performance metrics} (June 2026): the compute DP2 actually
      used.
\item \citeds{LDM-141}: catalogue row sizes and the object, source and
      forced-source counts per data release.
\item The USDF Qserv administrative dashboard (September 2026): the measured
      size of the DP2 catalogue.
\item \emph{2026 Hardware Initial Capacity Capture} (August 2026 revision):
      unit costs and the FY2026 purchase requests.
\item \emph{Rubin Hardware Summary for the March 2025 Joint Technical
      Meeting}: the purchase history.
\item RFC-1134: Prompt Processing product tiers and retention.
\end{itemize}

\subsection{Units}

Compute is given three ways. \emph{Core-hours} are what the estimates measure. \emph{Node-days} are how
campaigns are reported, at the \smCoresNode{}-core reference node.
\emph{Peak concurrent cores} across the processing window are what has to be
bought. Storage uses decimal units (1\,PB = 1000\,TB). Costs are in US dollars
of the year of purchase.
